% Generated draft: requires mathematical and source review.
\documentclass[11pt]{article}
\usepackage[margin=25mm]{geometry}
\usepackage{fontspec}
\setmainfont{Latin Modern Roman}
\usepackage{amsmath,amssymb,mathtools,unicode-math}
\setmathfont{Latin Modern Math}
\usepackage{enumitem,needspace,xurl,hyperref}
\setlength{\parindent}{0pt}
\setlength{\parskip}{5pt}
\setlength{\emergencystretch}{3em}
\newcommand{\N}{\mathbb{N}}
\newcommand{\Z}{\mathbb{Z}}
\newcommand{\Q}{\mathbb{Q}}
\newcommand{\R}{\mathbb{R}}
\newcommand{\C}{\mathbb{C}}
\newcommand{\F}{\mathbb{F}}
\newcommand{\A}{\mathbb{A}}
\newcommand{\PP}{\mathbb{P}}
\newcommand{\E}{\mathbb{E}}
\newcommand{\eps}{\varepsilon}
\newcommand{\dd}{\,\mathrm{d}}
\begin{document}
\section*{UnsolvedMath ID: 30001726}
\textbf{Source problem code:} catalogue entry
\subsection{Definitions and mathematical statement}
The supplied record does not define a mathematical table, a numerical table, or the temporal scope of the digital era. These terms are the subjects of a historiographical investigation; no formal class of tables, mathematical domain, or criterion for their disappearance or transformation is specified.

The research question attributed to the 2011 Oberwolfach mini-workshop is: What is a mathematical table in the era of digital computing? Did the development of digital programmable computers end the numerical table, or instead transform its algorithms and uses?

This is a historical and conceptual research question, rather than a specified mathematical proposition with explicit hypotheses and quantifiers. Its formulation leaves unresolved what counts as a table and what evidence would establish the end or transformation of numerical tables. A precise mathematical assertion cannot be extracted without introducing definitions and criteria absent from the supplied source record.
\subsection{Short English statement}
Investigate what mathematical tables became with digital computing, including whether programmable computers ended numerical tables or transformed their algorithms and uses. The supplied question is historiographical; its key terms and criteria require clarification before any mathematical formalization.
\subsection{Sources}
\begin{itemize}
\item Local UnsolvedMath registry, permanent ID \(30001726\), title \emph{Mathematical Tables in the Digital Era}, source problem code \texttt{catalogue entry}, in the Oberwolfach Reports Open Problems collection. The supplied record gives 2011 as the proposed year and identifies no individual proposer.
\item Supplied source citation: \emph{Mini-Workshop: History of Numerical and Graphical Tables} (2011), Oberwolfach Reports, report 12/2011, pp.\ 639--689, DOI \texttt{10.4171/owr/2011/12}. \url{https://doi.org/10.4171/owr/2011/12}. The supplied background describes Oberwolfach Reports as documenting mathematical workshops and their open-problem discussions.
\item Supplied literature source: \emph{Oberwolfach Report 12/2011}, \url{https://doi.org/10.4171/owr/2011/12}. Literature checked at: \texttt{2026-08-21}. This is supplied dated curation, not an independently verified literature assessment. It labels the question open and characterizes it as a continuing historical research question rather than a theorem with a definitive yes-or-no resolution; it requests independent verification of time-sensitive status and cited claims.
\item The supplied statement status is \texttt{corrected\_verified}. Its verification note says that the official report confirms the historiographical research question and that an extraction label was removed. These are supplied verification claims; the report has not been independently consulted for this draft.
\end{itemize}
\end{document}
